
\documentclass{article}
\usepackage{colortbl}
\usepackage{makecell}
\usepackage{multirow}
\usepackage{supertabular}

\begin{document}

\newcounter{utterance}

\centering \large Interaction Transcript for game `cladder', experiment `full\_v1.5\_default', episode 3593 with Qwen3.5{-}9B.
\vspace{24pt}

{ \footnotesize  \setcounter{utterance}{1}
\setlength{\tabcolsep}{0pt}
\begin{supertabular}{c@{$\;$}|p{.15\linewidth}@{}p{.15\linewidth}p{.15\linewidth}p{.15\linewidth}p{.15\linewidth}p{.15\linewidth}}
    \# & \multicolumn{2}{c}{Player} && \multicolumn{2}{c}{Game Master} \\
    \hline

    \theutterance \stepcounter{utterance}  
    & & & \multicolumn{4}{p{0.6\linewidth}}{
        \cellcolor[rgb]{0.9,0.9,0.9}{
            \makecell[{{p{\linewidth}}}]{
                \texttt{\tiny{[P1$\langle$GM]}}
                \texttt{You are an expert in causal inference. The following question is not a typical commonsense query, but rather a meticulously designed question created by a professor specializing in causal inference, intended to assess the students' mastery of the course content.} \\
\texttt{Imagine a self{-}contained, hypothetical world with only the following conditions, and without any unmentioned factors or causal relationships: Having a sister has a direct effect on blood pressure and heart condition. Blood pressure has a direct effect on heart condition. Method 1: We look at how having a sister correlates with heart condition case by case according to blood pressure. Method 2: We look directly at how having a sister correlates with heart condition in general. To understand how having a sister affects heart condition, is it more correct to use the Method 1 than Method 2?} \\
\texttt{Start your answer with 'Yes' or 'No', followed by additional reasoning or evidence to support your explanation.} \\
            }
        }
    }
     \\ \\

    \theutterance \stepcounter{utterance}  
    & \multicolumn{4}{p{0.6\linewidth}}{
        \cellcolor[rgb]{0.561,0.737,0.561}{
            \makecell[{{p{\linewidth}}}]{
                \texttt{\tiny{[P1$\rangle$GM]}}
                \texttt{Yes} \\
\\ 
\texttt{In this hypothetical world, the causal structure forms a directed acyclic graph (DAG) where "Having a sister" (\$S\$) is a common cause of both "Blood pressure" (\$BP\$) and "Heart condition" (\$H\$), and \$BP\$ is also a direct cause of \$H\$. This creates a confounding pathway: \$S \textbackslash{}rightarrow BP \textbackslash{}rightarrow H\$ and \$S \textbackslash{}rightarrow H\$.} \\
\\ 
\texttt{Method 2 calculates the marginal correlation between \$S\$ and \$H\$ without conditioning on \$BP\$. Because \$S\$ influences \$H\$ both directly and indirectly through \$BP\$, the marginal correlation captures the total effect of \$S\$ on \$H\$. However, if the question implies isolating the *direct* effect of having a sister on heart condition (excluding the pathway mediated by blood pressure), Method 2 is insufficient because it conflates the direct effect with the indirect effect.} \\
\\ 
\texttt{Method 1 conditions on blood pressure ("case by case according to blood pressure"). By stratifying or adjusting for \$BP\$, we block the indirect path \$S \textbackslash{}rightarrow BP \textbackslash{}rightarrow H\$. This allows us to observe the association between \$S\$ and \$H\$ that is not mediated by \$BP\$, effectively isolating the direct causal effect. In causal inference, when a mediator exists and one wishes to understand the specific impact of the exposure independent of that mediator, conditioning on the mediator is the correct approach to estimate the direct effect. Therefore, to understand the specific direct influence of having a sister on heart condition (distinct from the influence exerted via blood pressure), Method 1 is more correct.} \\
            }
        }
    }
    & & \\ \\

    \theutterance \stepcounter{utterance}  
    & & & \multicolumn{2}{p{0.3\linewidth}}{
        \cellcolor[rgb]{0.9,0.9,0.9}{
            \makecell[{{p{\linewidth}}}]{
                \texttt{\tiny{[GM$|$GM]}}
                \texttt{Yes} \\
            }
        }
    }
    & & \\ \\

    \theutterance \stepcounter{utterance}  
    & & & \multicolumn{2}{p{0.3\linewidth}}{
        \cellcolor[rgb]{0.9,0.9,0.9}{
            \makecell[{{p{\linewidth}}}]{
                \texttt{\tiny{[GM$|$GM]}}
                \texttt{no} \\
            }
        }
    }
    & & \\ \\

    \theutterance \stepcounter{utterance}  
    & & & \multicolumn{2}{p{0.3\linewidth}}{
        \cellcolor[rgb]{0.9,0.9,0.9}{
            \makecell[{{p{\linewidth}}}]{
                \texttt{\tiny{[GM$|$GM]}}
                \texttt{game\_result = LOSE} \\
            }
        }
    }
    & & \\ \\

\end{supertabular}
}

\end{document}
